\newpage
\subsection{Тестирование}

Тестирование интеллектуальной системы анализа обитаемости экзопланет направлено на проверку корректности работы всех компонентов — от валидации входных параметров до визуализации результатов анализа. Ниже представлены этапы тестирования с иллюстрациями, демонстрирующими работу системы в реальных условиях.

\subsubsection{Тестирование валидации входных параметров}

Ключевым аспектом работы системы является корректная валидация параметров экзопланеты. В процессе тестирования проверялась работа следующих компонентов:

\begin{itemize}
    \item валидация числовых значений радиуса, температуры и плотности планеты;
    \item обработка некорректных значений (отрицательные числа, значения вне допустимого диапазона);
    \item работа подсказок и вспомогательной информации при вводе данных.
\end{itemize}

На рисунке 2.1 представлена форма ввода параметров с примером заполнения данных для экзопланеты:

[Здесь будет скриншот формы ввода параметров с подсказками]

При тестировании особое внимание уделялось проверке граничных значений параметров:
\begin{itemize}
    \item Радиус: допустимый диапазон от 0.1 до 3.0 радиусов Земли
    \item Температура: от 50К до 1000К
    \item Плотность: от 0.01 до 3.0 относительно плотности Земли
\end{itemize}

\subsubsection{Тестирование процесса анализа}

После валидации входных данных система выполняет комплексный анализ потенциальной обитаемости экзопланеты. Тестирование этого этапа включало проверку:

\begin{itemize}
    \item корректности расчета индекса подобия Земле (ESI);
    \item точности определения вероятности наличия жидкой воды;
    \item правильности анализа атмосферных параметров;
    \item генерации рекомендаций на основе рассчитанных показателей.
\end{itemize}

На рисунке 2.2 показан процесс анализа с индикацией прогресса выполнения каждого этапа:

[Здесь будет скриншот процесса анализа с прогресс-баром]

\subsubsection{Тестирование визуализации результатов}

Особое внимание было уделено тестированию компонентов визуализации результатов:

\begin{itemize}
    \item корректность построения 3D-модели планеты с учетом введенных параметров;
    \item точность отображения графиков и диаграмм сравнения с Землей;
    \item работа интерактивных элементов управления 3D-моделью;
    \item корректность отображения индикаторов обитаемости.
\end{itemize}

На рисунке 2.3 представлены результаты анализа с различными видами визуализации:

[Здесь будет скриншот страницы результатов с графиками и 3D-моделью]

\subsubsection{Тестирование сравнительного анализа}

Система позволяет сравнивать параметры различных экзопланет. При тестировании этого функционала проверялись:

\begin{itemize}
    \item корректность поиска похожих экзопланет в базе данных;
    \item точность расчета сравнительных характеристик;
    \item наглядность представления результатов сравнения.
\end{itemize}

На рисунке 2.4 показан пример сравнительного анализа нескольких экзопланет:

[Здесь будет скриншот сравнительного анализа]

\subsubsection{Результаты тестирования}

В ходе тестирования были проверены все основные функции системы и подтверждена их корректная работа. Особое внимание было уделено точности математических расчетов и качеству визуализации данных. Тестирование показало, что система:

\begin{itemize}
    \item корректно валидирует входные параметры экзопланет;
    \item точно выполняет расчеты индекса подобия Земле и других показателей;
    \item качественно визуализирует результаты анализа;
    \item обеспечивает удобный и информативный интерфейс.
\end{itemize}

Все выявленные в процессе тестирования недочеты были успешно устранены, что позволяет говорить о готовности системы к практическому использованию в области анализа потенциально обитаемых экзопланет. 